\answer{ex:01:1}{(a)~$8\cdot10^{10}$ \nt{GLUT} --- $10$ lần dân số Trái Đất; nơ-ron quảng bá $\approx1.9\cdot10^6$ --- một thành phố cỡ Viên. (b)~$\approx8\cdot10^{14}$ đầu tận cùng, trong đó của nơ-ron quảng bá $\approx3\cdot10^{11}$, tỷ lệ $\approx4\cdot10^{-4}$ --- lớn hơn $\approx16$ lần tỷ lệ các nơ-ron này ($2.2\cdot10^{-5}$).}
\answer{ex:01:2}{(a)~$\tau_c\ln(c_0/c_*)\approx4.6\ms$. (b)~Đường truyền rảnh khi $T\ge\tau_c\ln101$: tới $\approx22$ lần nhả mỗi giây thay vì $\approx220$.}
\answer{ex:01:3}{$V(s_k)=\gamma^{K-1-k}r$; $\delta_{\text{đèn}}=\gamma^Kr=re^{-T/\tau_\gamma}\approx0.60\,r$ với mọi $\Delta$, $\tau_\gamma\approx1.95\,\text{s}$.}
\answer{ex:01:4}{$n^*=93$, $47$, $25$, $12$: $n^*\approx5/\alpha$ khi $\alpha$ nhỏ.}
\answer{ex:01:5}{\nt{GLUT}: $1\to10\to10^4\to10^7$ --- $\approx10^7$ nơ-ron, $4$ mắt xích, $4\ms$. \nt{DOPA}: $1\to10\to3.2\cdot10^5$ --- $\approx3\cdot10^5$ nơ-ron, $3$ mắt xích, ít hơn $30$ lần; cùng bậc với số nơ-ron \nt{DOPA}.}
\answer{ex:01:6}{$\lambda=f_EJ_E-f_IJ_I$, suy ra $J_I=37.5$; tỷ số dòng $f_IJ_I/f_EJ_E=0.94$; thác hoạt động ($\lambda\ge1$) xảy ra khi ức chế chỉ yếu đi $6.7\,\%$.}
\answer{ex:01:7}{$n\ge57$ lần trình bày cho mỗi độ trễ trong mỗi điều kiện. Cùng sự suy giảm đó có thể đến, chẳng hạn, từ sai số của ước lượng thời gian bên trong, tăng theo $T$.}
